% =====================================================================
\chapter{Conclusion and Future Work}
\label{ch:conclusion}
% =====================================================================

\section{Summary}
\label{sec:conclusion:summary}

This thesis set out to quantify and reduce the temporal gap between
firmware running on physical hardware and the same firmware rehosted in
the Renode emulation framework.

On the \emph{quantification} side (RQ1), a twin measurement
methodology---bit-identical firmware, GPIO-edge observables, a common
statistical pipeline---revealed that the gap is not a simple slowdown
or speedup. Baseline Renode reproduces \emph{mean} timing to within
fractions of a percent, yet distorts timing \emph{variance}
bidirectionally: tick-quantised sleeps appear $5\times$ too
deterministic, while timer callbacks and ISR latencies appear up to
$23\times$ too noisy, an artefact of the \SI{100}{\micro\second}
scheduling quantum. A taxonomy of five deviation sources
(\srcS{1}--\srcS{5}) attributes each effect to a specific emulation
mechanism.

On the \emph{mitigation} side (RQ2), the Temporal Fidelity Layer
demonstrates that a substantial share of fidelity can be recovered
without modifying either the emulator or the firmware: three
runtime-loaded C\# peripherals inject calibrated clock drift and
jitter, monitor scheduler wakeup deviation, and model inter-node
propagation delay, all configured declaratively in platform files.
Across seven statistically evaluated scenarios, TFL reduced the
geometric-mean jitter-ratio fold error from $7.6\times$ to $2.9\times$
and brought four scenarios within a factor of $1.7$ of hardware
variance. The evaluation also exposed a miscalibration (proportional
rather than absolute sleep jitter) that only hardware-grounded
measurement could reveal---underlining the thesis's central
methodological claim: emulator fidelity must be \emph{measured}, not
assumed.

\section{Future Work}
\label{sec:conclusion:future}

\subsection{Model Calibration from Traces}
\label{sec:future:calibration}

The T3 regression shows that hand-parameterised models have limits. A
natural next step is closing the loop: derive TFL parameters
automatically from a short hardware characterisation run (e.g.\ fit an
absolute uniform tick-rounding term plus a load-dependent Gaussian
term per primitive), turning TFL into a self-calibrating fidelity
layer.

\subsection{Addressing S1 and S2}
\label{sec:future:s1s2}

Instruction-timing abstraction and interrupt recognition granularity
remain untouched by TFL. Promising directions include per-block timing
annotations derived from cycle-accurate reference runs, and
finer-grained interrupt scheduling below the quantum---ideally
upstreamed into Renode's core rather than layered on top.

\subsection{Generalisation}
\label{sec:future:generalisation}

The methodology transfers directly to other targets (Cortex-M4/M7 with
caches, RISC-V), other RTOSes (FreeRTOS, RIOT), and multi-node setups
where \srcS{5} dominates. Multi-node experiments with
\texttt{TFL\_BusDelayModel} on realistic protocols (CAN, BLE mesh)
would test whether distribution-level fidelity suffices to expose real
protocol race conditions in emulation.

\subsection{Fidelity as a CI Metric}
\label{sec:future:ci}

Finally, the jitter-ratio metric and the deadline-monitor CSV stream
are cheap enough to run in continuous integration. Tracking $J$ per
release would let projects detect fidelity regressions in their
emulated test environments the same way they track code coverage
today.
